% GENERATED FILE. Edit canonical lessons, then run npm run generate:books.
% canonical-source-hash: fnv1a64:3d36964f2382ce66
% canonical-lessons: UR-C426-khalis, UR-C426-asli, UR-C426-milavat, UR-C426-ragarna, UR-C426-chamkana

\chapter{Pure and Genuine}
\label{ch:ur-pure-and-genuine}
\hlchaptermodality{426}

\begin{chapteropening}
\textbf{By the end of this chapter:} \emph{I can ask whether goods are pure and genuine.}

Complete the last lesson of chapter 426: I can ask whether goods are pure and genuine.
\end{chapteropening}

\section[\texorpdfstring{\ur{خالص} (khālis)}{خالص (khālis)}]{\ur{خالص} (khālis) — unmixed; pure}

\label{lesson:UR-C426-khalis}



\begin{quote}
Before the new one: say the Urdu for loose change, then the Urdu for to hand (to someone).
\end{quote}

\subsection*{You'll want to know: \ur{خالص}}
\textbf{\ur{خالص}} — \emph{khālis} — \textquotedblleft{}unmixed; pure\textquotedblright{}.

\begin{quote}
\ur{یہ} \ur{دودھ} \ur{خالص} \ur{ہے۔}
\end{quote}

\begin{quote}
\emph{yih dūdh khālis hai.}
\end{quote}

\begin{quote}
This milk is pure.
\end{quote}

\begin{cousinweb}
\textbf{\ur{خالص}} (\emph{khālis}) is usually traced, through Persian, to Arabic \emph{khāliṣ}, \textquotedblleft{}pure, unmixed\textquotedblright{}.
\end{cousinweb}

\begin{grammarlens}[title={What you've built}]
One more word for getting around.
\end{grammarlens}

\subsection*{Your turn}
\begin{itemize}
\raggedright
  \item \emph{Say it:} \emph{khālis}
  \item \emph{Say it:} \emph{khālis}, once more
  \item \emph{Say it:} \emph{yih dūdh khālis hai.} — This milk is pure.
  \item \emph{Recall:} say the Urdu for loose change, then the Urdu for to hand (to someone), then say \emph{khālis} again
\end{itemize}

\begin{tcolorbox}[breakable,colback=teal!4,colframe=teal!35!black,title={Before you move on}]
What does \ur{خالص} mean? (\textquotedblleft{}Unmixed; pure\textquotedblright{}.) Say it once more.
\end{tcolorbox}
\section[\texorpdfstring{\ur{اصلی} (aslī)}{اصلی (aslī)}]{\ur{اصلی} (aslī) — authentic; real}

\label{lesson:UR-C426-asli}



\begin{quote}
Before the new one: say the Urdu for to hand (to someone), then the Urdu for unmixed.
\end{quote}

\subsection*{You'll want to know: \ur{اصلی}}
\textbf{\ur{اصلی}} — \emph{aslī} — \textquotedblleft{}authentic; real\textquotedblright{}.

\begin{quote}
\ur{یہ} \ur{اصلی} \ur{سونا} \ur{ہے،} \ur{نقلی} \ur{نہیں۔}
\end{quote}

\begin{quote}
\emph{yih aslī sonā hai, naqlī nahīṅ.}
\end{quote}

\begin{quote}
This is real gold, not fake.
\end{quote}

\begin{cousinweb}
\textbf{\ur{اصلی}} (\emph{aslī}) is built on \emph{asl}, \textquotedblleft{}the origin, the real thing\textquotedblright{}, usually traced, through Persian, to Arabic \emph{aṣl}, \textquotedblleft{}a root\textquotedblright{}, with the ending \emph{-ī}.
\end{cousinweb}

\begin{grammarlens}[title={What you've built}]
One more word for getting around.
\end{grammarlens}

\subsection*{Your turn}
\begin{itemize}
\raggedright
  \item \emph{Say it:} \emph{aslī}
  \item \emph{Say it:} \emph{aslī}, once more
  \item \emph{Say it:} \emph{yih aslī sonā hai, naqlī nahīṅ.} — This is real gold, not fake.
  \item \emph{Recall:} say the Urdu for to hand (to someone), then the Urdu for unmixed, then say \emph{aslī} again
\end{itemize}

\begin{tcolorbox}[breakable,colback=teal!4,colframe=teal!35!black,title={Before you move on}]
What does \ur{اصلی} mean? (\textquotedblleft{}Authentic; real\textquotedblright{}.) Say it once more.
\end{tcolorbox}
\section[\texorpdfstring{\ur{ملاوٹ} (milāvaṭ)}{ملاوٹ (milāvaṭ)}]{\ur{ملاوٹ} (milāvaṭ) — adulteration; mixing in}

\label{lesson:UR-C426-milavat}



\begin{quote}
Before the new one: say the Urdu for unmixed, then the Urdu for authentic.
\end{quote}

\subsection*{You'll want to know: \ur{ملاوٹ}}
\textbf{\ur{ملاوٹ}} — \emph{milāvaṭ} — \textquotedblleft{}adulteration; mixing in\textquotedblright{}.

\begin{quote}
\ur{اس} \ur{دودھ} \ur{میں} \ur{پانی} \ur{کی} \ur{ملاوٹ} \ur{ہے۔}
\end{quote}

\begin{quote}
\emph{is dūdh meṅ pānī kī milāvaṭ hai.}
\end{quote}

\begin{quote}
This milk has been watered down.
\end{quote}

Feminine.

\begin{cousinweb}
\textbf{\ur{ملاوٹ}} (\emph{milāvaṭ}) is built on \textbf{\ur{ملانا}}, \textquotedblleft{}to mix\textquotedblright{}, from Sanskrit \emph{milati}, \textquotedblleft{}meets\textquotedblright{} (CDIAL 10133), with the noun ending \emph{-āvaṭ}: in a shop it means something cheaper mixed in.
\end{cousinweb}

\begin{grammarlens}[title={What you've built}]
One more word for getting around.
\end{grammarlens}

\subsection*{Your turn}
\begin{itemize}
\raggedright
  \item \emph{Say it:} \emph{milāvaṭ}
  \item \emph{Say it:} \emph{milāvaṭ}, once more
  \item \emph{Say it:} \emph{is dūdh meṅ pānī kī milāvaṭ hai.} — This milk has been watered down.
  \item \emph{Recall:} say the Urdu for unmixed, then the Urdu for authentic, then say \emph{milāvaṭ} again
\end{itemize}

\begin{tcolorbox}[breakable,colback=teal!4,colframe=teal!35!black,title={Before you move on}]
What does \ur{ملاوٹ} mean? (\textquotedblleft{}Adulteration; mixing in\textquotedblright{}.) Say it once more.
\end{tcolorbox}
\section[\texorpdfstring{\ur{رگڑنا} (ragaṛnā)}{رگڑنا (ragaṛnā)}]{\ur{رگڑنا} (ragaṛnā) — to rub; to scrub}

\label{lesson:UR-C426-ragarna}



\begin{quote}
Before the new one: say the Urdu for authentic, then the Urdu for adulteration.
\end{quote}

\subsection*{You'll want to know: \ur{رگڑنا}}
\textbf{\ur{رگڑنا}} — \emph{ragaṛnā} — \textquotedblleft{}to rub; to scrub\textquotedblright{}.

\begin{quote}
\ur{سنار} \ur{نے} \ur{سونا} \ur{پتھر} \ur{پر} \ur{رگڑا۔}
\end{quote}

\begin{quote}
\emph{sunār ne sonā patthar par ragaṛā.}
\end{quote}

\begin{quote}
The goldsmith rubbed the gold on a stone.
\end{quote}

Past, with \ur{نے} (he or she …): \ur{اس} \ur{نے} \ur{برتن} \ur{رگڑے}, \emph{us ne bartan ragaṛe}.

\begin{cousinweb}
\textbf{\ur{رگڑنا}} (\emph{ragaṛnā}) is inherited: CDIAL lists the Hindi cousin \emph{ragaṛnā}, \textquotedblleft{}to rub\textquotedblright{}, under \emph{ragg}, a form it reconstructs (CDIAL 10558).
\end{cousinweb}

\begin{grammarlens}[title={What you've built}]
One more word for getting around.
\end{grammarlens}

\subsection*{Your turn}
\begin{itemize}
\raggedright
  \item \emph{Say it:} \emph{ragaṛnā}
  \item \emph{Say it:} \emph{ragaṛnā}, once more
  \item \emph{Say it:} \emph{sunār ne sonā patthar par ragaṛā.} — The goldsmith rubbed the gold on a stone.
  \item \emph{Recall:} say the Urdu for authentic, then the Urdu for adulteration, then say \emph{ragaṛnā} again
\end{itemize}

\begin{tcolorbox}[breakable,colback=teal!4,colframe=teal!35!black,title={Before you move on}]
What does \ur{رگڑنا} mean? (\textquotedblleft{}To rub; to scrub\textquotedblright{}.) Say it once more.
\end{tcolorbox}
\section[\texorpdfstring{\ur{چمکانا} (chamkānā)}{چمکانا (chamkānā)}]{\ur{چمکانا} (chamkānā) — to polish; to make (something) shine}

\label{lesson:UR-C426-chamkana}



\begin{quote}
Before the new one: say the Urdu for adulteration, then the Urdu for to rub.
\end{quote}

\subsection*{You'll want to know: \ur{چمکانا}}
\textbf{\ur{چمکانا}} — \emph{chamkānā} — \textquotedblleft{}to polish; to make (something) shine\textquotedblright{}.

\begin{quote}
\ur{اس} \ur{نے} \ur{کپڑے} \ur{سے} \ur{انگوٹھی} \ur{چمکائی۔}
\end{quote}

\begin{quote}
\emph{us ne kapṛe se angūṭhī chamkāī.}
\end{quote}

\begin{quote}
She polished the ring with a cloth.
\end{quote}

Past, with \ur{نے} (he or she …): \ur{اس} \ur{نے} \ur{جوتے} \ur{چمکائے}, \emph{us ne jūte chamkāe}.

\begin{cousinweb}
\textbf{\ur{چمکانا}} (\emph{chamkānā}) is built on \emph{chamaknā}, \textquotedblleft{}to shine\textquotedblright{}, which CDIAL lists under \emph{camakka}, a form it reconstructs (CDIAL 4676): this is the form for making something shine.
\end{cousinweb}

\begin{grammarlens}[title={What you've built}]
One more word for getting around.
\end{grammarlens}

\subsection*{Your turn}
\begin{itemize}
\raggedright
  \item \emph{Say it:} \emph{chamkānā}
  \item \emph{Say it:} \emph{chamkānā}, once more
  \item \emph{Say it:} \emph{us ne kapṛe se angūṭhī chamkāī.} — She polished the ring with a cloth.
  \item \emph{Recall:} say the Urdu for adulteration, then the Urdu for to rub, then say \emph{chamkānā} again
\end{itemize}

\begin{tcolorbox}[breakable,colback=teal!4,colframe=teal!35!black,title={Before you move on}]
What does \ur{چمکانا} mean? (\textquotedblleft{}To polish; to make (something) shine\textquotedblright{}.) Say it once more.
\end{tcolorbox}
